\documentclass[11pt,a4paper]{article}

\usepackage[margin=2.3cm]{geometry}
\usepackage[T1]{fontenc}
\usepackage[utf8]{inputenc}
\usepackage{lmodern}
\usepackage{microtype}
\usepackage{setspace}
\usepackage{booktabs}
\usepackage{longtable}
\usepackage{array}
\usepackage{enumitem}
\usepackage{xcolor}
\usepackage{hyperref}

\onehalfspacing
\setlist[itemize]{leftmargin=1.3em,itemsep=0.3em,topsep=0.35em}
\setlength{\parskip}{0.4em}
\setlength{\parindent}{0pt}
\hypersetup{colorlinks=true,linkcolor=black,urlcolor=blue}

\title{\vspace{-1.2em}What was wrong in Kerry's specification,\\[0.25em]
\large and what was changed so it matches the code}
\author{ECO5040W Group 3}
\date{27 September 2026}

\begin{document}
\maketitle
\vspace{-1.4em}

This note is about \texttt{ECO5040W\_Business\_And\_Tech\_Specification\_kerrys.docx} only.
It records what was already right, what was false when we first read it, and the wording that replaced the false parts.
The file still has Kerry's 12 tables and all six figures.
Nothing in those figures was redrawn.

\section{Are the tables, figures, and claims good?}

The tables are good where they state a number or a rule the code actually uses.
The competitor percentages were checked and left alone.
The limit tiers are the brief's R0 / R0 and R3{,}000 / R25{,}000.
The fee example for R1{,}000 at the fallback rate matches the quote code.
Several \emph{Implemented in} cells were too neat: they described a stricter or simpler system than the one that runs.
Those cells were rewritten.
The requirement column was left as the brief wrote it, except FR-12b and FR-12c, whose texts had been swapped, and FR-19, which now says an administrator confirms cash-in because that is the only route in the code.

The figures are Kerry's, and they are still the right kind of figure: three journey diagrams, an architecture diagram, an entity diagram, and a picture of the configuration tables.
They were not regenerated.
One diagram sits under the wrong heading, and the entity diagram is behind the schema.
Those are noted below and were not ``fixed'' by drawing a new picture.

The claims are good now.
The sentences that said the prototype was an anti-money-laundering system, enforced the R2~million annual allowance, showed an exact cash-out amount, kept client money in a segregated bank account, or was cheaper than WorldRemit have been replaced.
What is still rough is draft clutter on the surface of the file (an unfinished heading, a missing space, comments).
That clutter is not a wrong claim about the system.

\section{Left alone, because it was already right}

\begin{itemize}
\item The World Bank corridor figures in the opening and in the competitor table: MoneyGram 15.43\%, Western Union 31.81\%, Mukuru 10.28\%, WorldRemit 5.15\%, corridor average 12.49\%.
\item The default fees: R25 fixed, 1\% of the amount, 2\% exchange-rate margin, 1.5\% cash-out fee, stored as the fractions 0.0100, 0.0200, and 0.0150.
\item The quote algebra: fee first, then the margin on the mid-market rate, then the token amount.
\item The limit amounts in the tier table, once the header spelling was corrected.
\item The six images. Their bytes were checked before and after every save.
\end{itemize}

\section{Claims that were false}

Each row is the claim as it stood, then the replacement that matches the code.

\begin{description}[leftmargin=0pt,style=nextline,font=\normalfont\bfseries]

\item[Cheaper than the digital providers.]
The closing paragraph of the comparison said the platform targets a total cost below the digital providers by settling in a stablecoin.
At the default fees that is not true, and it is not under the 3\% target.
The paragraph now states the R25 fee, the 1\%, the 2\% margin, and the estimated 1.5\% cash-out fee, and it says that combination is not below those providers and not under 3\%.
The WorldRemit and incumbent percentages in the same paragraph were not touched.

\item[Anti-money-laundering controls.]
The regulatory section said the KYC process and the daily and monthly limits are the transaction controls a monitored service requires.
The prototype collects mock KYC, blocks sending and cashing out until an administrator approves, and enforces rand caps.
Those caps are not a monitoring or reporting engine.
The section now says that.

\item[The R2 million allowance.]
The text said the platform treats RLUSD transfers like any other transfer and applies the same limits, after describing the Reserve Bank's annual allowance.
The prototype does not enforce that allowance.
It enforces configurable daily and monthly caps, R3{,}000 and R25{,}000 for a verified sender.
The sentence now says a live dealer would have to apply the annual allowance, and this prototype does not.

\item[The exact amount the recipient will receive.]
Senders were said to see the exact amount before they confirm.
The quote's cash-out fee and payout are estimates.
Cash-out recalculates them later.
The sentence now says ``estimate'', and that the payout is recalculated when the recipient cashes out.

\item[Client money kept apart from the business.]
Balances were said to be recorded separately so a customer's holding is never mixed with the money the business uses to operate.
Spendable balance is an internal ledger.
One platform treasury pays the recipient's custodial account, and the platform holds that key.
Cash-in and fiat payout are simulated.
There is no separate bank account for client money.
The paragraph now says that.

\item[Private keys decrypted only by the signer.]
Both the custody paragraph and the security section said the secret is decrypted only by the code that signs.
At the time, loading the wallet row decrypted it.
The code was then changed: the row loads ciphertext, and the seed is decrypted only when settlement, cash-out completion, or a TrustSet is about to sign.
The specification was updated a second time to that behaviour.
Wallet responses still omit the secret.
The KYC identification number is still decrypted when the KYC record is read, which is intentional.

\item[Rounding.]
``Fees are rounded to the nearest cent and token amounts to six decimal places'' was too coarse.
The rand transaction fee and the fiat payout use two decimal places.
The cash-out fee and the token amounts use six, halves rounded up.
A worked R1{,}000 at the 18.50 fallback was added and checked: fee R35.00, effective rate 18.87, net R965.00, 51.139375 tokens, estimated cash-out fee 0.767091, estimated payout 50.372284 tokens, not rand.

\item[Who the limits apply to, and what R0 does.]
The line ``whether is verified or unverified'' was incomplete, and the next sentence said the zero limits exist so unverified users cannot send.
Sending is blocked earlier: \texttt{POST /remittances} returns 403 until KYC is approved, so an unverified sender never reaches the R0 check.
R0 is the configured tier shown on \texttt{GET /limits/me}.

\item[What counts toward the cap.]
Usage was easy to read as ``everything until the day or month ends''.
The code sums gross \texttt{zar\_amount}.
A quote is refused only when usage plus the new amount is strictly greater than the cap, so a send that lands on the cap is allowed.
Cancel and expiry release the amount before the period ends.
\texttt{settlement\_failed} still counts.
Two quotes sent together are not locked, so both can pass.
The limits section now says each of those.

\end{description}

\section{Tables}

\subsection{Headers and one status that was not a status}

\begin{itemize}
\item The limit table header said ``Status'' and ``Montly limit''.
It now says Tier, Daily limit, Monthly limit.
The two rows were already the right amounts.
\item The status table had a row ``quoted (past TTL)''.
That is not a stored status.
The row remains \texttt{quoted}.
The cell now says ``quoted, after expiry'', and the next cell says cash-in returns 409 once \texttt{expires\_at} has passed and the amount no longer counts.
\end{itemize}

\subsection{FR-12b and FR-12c were swapped}

The course brief's FR-12b is ``provision a custodial wallet only once the beneficiary is linked''.
FR-12c is ``if nobody matches, leave the beneficiary unlinked and do not settle''.
Kerry's table had those two requirement texts the other way around.
The requirement text and the implementation text were swapped together so the identifiers match the brief.

The implementation text was then made specific.
FR-12b: a wallet row is created when the beneficiary is linked; the Testnet account is created at the first settlement, and only for that linked user.
FR-12c: the record stays unlinked; if cash-in is confirmed anyway, the worker refuses, credits nothing, and marks the transfer \texttt{settlement\_failed}.

\subsection{Implemented-in cells that oversimplified the code}

{\small
\begin{longtable}{@{}>{\raggedright\arraybackslash}p{1.35cm}>{\raggedright\arraybackslash}p{6.3cm}>{\raggedright\arraybackslash}p{7.1cm}@{}}
\toprule
ID & What the cell used to imply & What it says now \\
\midrule
\endfirsthead
\toprule
ID & What the cell used to imply & What it says now \\
\midrule
\endhead
\bottomrule
\endfoot
FR-06 &
Resubmission overwrites the application until it is approved. &
The API overwrites a pending or rejected application and returns 409 once it is approved.
The screen locks the form while it is pending. \\
FR-08 &
Reject ``with a reason'', as if the API required one. &
The API returns 422 if the reason is missing or blank.
The admin screen asks for one before sending.
The code was changed so this is true. \\
FR-09 &
The send screen shows ``KYC pending''. &
The heading is ``KYC not approved'', with a line for pending, rejected, or not submitted.
The API returns 403. \\
FR-10 &
All five beneficiary fields are stored, with no caveat. &
\texttt{payout\_currency} is stored and shown.
Quotes do not use it. \\
FR-12a &
Auto-match on add and on register. &
Mobile or email is enough.
The first matching account wins.
Editing the profile does not relink.
Email matching is case-insensitive; account email uniqueness is case-sensitive. \\
FR-16 &
Checked before insert; cancelled and expired quotes do not count. &
Checked on gross rand.
Refused only when usage plus the amount is strictly over the cap.
Two simultaneous quotes are not row-locked. \\
FR-19 &
The requirement said an administrator or a mock service. &
The requirement now says an administrator confirms the payment.
The implementation cell says there is no mock payment-provider route.
The brief's ``or'' is satisfied by the administrator. \\
FR-22 &
A worker reads the stream, which could be read as a process that stays up. &
The script or the admin action drains the stream once.
It does not stay running between calls. \\
FR-25 &
A duplicate message cannot credit twice, full stop. &
One message per remittance, a stored hash is not paid again, and a completed retry returns 409.
Before submitting, the worker looks up a validated payment whose memo is this remittance and adopts that hash.
Two overlapping runs can still both miss the lookup and pay twice. \\
FR-26 &
One commit after validation, with no mention of a crash. &
The commit stores the hash, credits the balance, and sets settled.
If the process dies before that commit, the next run adopts the payment by its memo and then credits the wallet. \\
FR-28 &
``Spendable and on-chain balances'', as if the on-chain figure were read from the ledger. &
The on-chain figure is spendable plus requested or approved cash-out amounts.
It is not a live trust-line read. \\
FR-28a &
The history page shows the quote breakdown, timeline, hash, and QR code. &
The history page lists amount, status, and tracking reference.
The breakdown, timeline, hash, and QR code are on the track page and the quote page. \\
FR-32 &
Complete burns the reserved tokens; fail refunds spendable. &
Complete burns the gross reserved amount.
The fiat payout is the net after the fee and is not paid by a bank.
Fail refunds that gross amount to the cash-out user. \\
FR-35 &
The admin screen approves, completes, or fails a cash-out. &
The same, plus: fail refunds the cash-out user, the recipient, not the remittance sender.
The admin screen's own sentence was wrong in the same way and was corrected in the code. \\
\bottomrule
\end{longtable}
}

\section{Figures}

The six images are still \texttt{image1.png} through \texttt{image6.png}.

\begin{itemize}
\item \texttt{image1} is the register, KYC, beneficiary, quote, and cancel sequence.
\item \texttt{image3} is cash-in and settlement.
\item \texttt{image2} is wallet, cash-out, and tracking.
It is also placed under Onboarding, which is the wrong heading for that picture.
Moving it would have meant editing the drawing anchors.
It was left where Kerry put it.
\item \texttt{image4} is the architecture diagram.
The prose under it was empty and was filled.
One sentence was added later: the prototype uses both wallet options in the brief, a separate Testnet account per linked recipient and an internal spendable ledger, and the on-chain figure is not a live trust-line balance.
\item \texttt{image5} is the entity diagram.
It does not show the platform wallet, the fee row, or the limit tiers, and its KYC status list is behind the code.
It was not redrawn.
The database section next to it names those tables and the three Alembic revisions.
\item \texttt{image6} is the configuration tables, including the fee defaults 25.00, 0.0100, 0.0200, and 0.0150.
Those match the code.
\end{itemize}

No figure was replaced, and no new diagram was generated.

\section{Code that the document had to follow}

Two small code changes were made because the document would otherwise have stayed wrong, or the code would have stayed looser than the brief.

\begin{itemize}
\item Rejecting a KYC application now requires a non-empty reason.
A missing or blank reason returns 422.
The FR-08 cell describes that.
\item The admin cash-out page said failing a cash-out refunds the sender.
It refunds the cash-out user.
The page and the FR-35 cell now say that.
\end{itemize}

A larger code change followed the FR-25 hole.
The worker now looks up an existing validated payment by the remittance memo before it submits another one.
The FR-25 cell, the settlement paragraph, the assumptions paragraph, and the FR-26 cell were updated to that behaviour, including the remaining case: two overlapping worker runs can still both pay.

Wallet secrets were changed in the same pass.
They load as ciphertext and are decrypted only at the signing call.
The two security sentences were updated so they no longer describe the old decrypt-on-load behaviour, and they do not claim the stricter rule until the code had it.

\end{document}
